% Resume source. Compile with pdflatex from this directory.
\documentclass[11pt,letterpaper]{article}
\usepackage[margin=0.65in]{geometry}
\usepackage[T1]{fontenc}
\usepackage[utf8]{inputenc}
\usepackage{helvet}
\renewcommand{\familydefault}{\sfdefault}
\usepackage{enumitem}
\usepackage{titlesec}
\usepackage{xcolor}
\usepackage[hidelinks]{hyperref}
\input{glyphtounicode}
\pdfgentounicode=1
\hypersetup{pdftitle={Haojie Ye - AI Systems and GPU Architecture},pdfauthor={Haojie Ye},pdfsubject={Professional Resume},pdfcreator={LaTeX}}
\urlstyle{same}
\pagestyle{empty}
\setlength{\parindent}{0pt}
\setlength{\parskip}{5pt}
\raggedright
\titleformat{\section}{\large\bfseries}{}{0pt}{}
\titlespacing*{\section}{0pt}{14pt}{6pt}
\setlist[itemize]{leftmargin=13pt,itemsep=4pt,topsep=3pt,parsep=0pt}
\newcommand{\role}[4]{\textbf{#1}\hfill #2\\{#3}\hfill {#4}\par}
\newcommand{\paper}[4]{\textbf{#1}\\#2\\\href{#4}{#3}\par\vspace{3pt}}
\begin{document}
\fontsize{10.5}{13}\selectfont

{\LARGE\textbf{Haojie Ye}}\\[4pt]
{\large AI Systems and GPU Architecture}\\[5pt]
\href{mailto:yehaojie@umich.edu}{yehaojie@umich.edu}\enspace\textbar\enspace
\href{https://linestro.github.io/}{linestro.github.io}\enspace\textbar\enspace
\href{https://www.linkedin.com/in/haojie-ye-91b935106/}{LinkedIn}\enspace\textbar\enspace
\href{https://scholar.google.com/citations?user=8Mk6B6gAAAAJ}{Google Scholar}

AI performance architect with a Ph.D. in Computer Science and Engineering and a research focus on computer architecture. Connects generative AI workloads to GPU design and product decisions through validated performance models, agent-assisted analysis tools, and investigation of compute, memory, and communication bottlenecks.

\section{Experience}
\role{NVIDIA}{Oct 2024 -- Present}{Deep Learning Performance Architect}{Santa Clara, CA}
\begin{itemize}
\item Led performance modeling studies for LLM inference and video generation, translating model assumptions and execution strategies into GPU bottleneck analysis and product tradeoff recommendations.
\item Implemented and validated simulation capabilities for evolving AI workloads, including sparse attention and parallel execution. Checked model shapes, execution mappings, and run assumptions to strengthen performance comparisons.
\item Independently developed an agent-assisted performance-attribution component that reconciles execution mappings, explains differences between projections, and surfaces outliers. Integrated it into a broader analysis workflow for repeatable comparison and targeted debugging.
\item Evaluated compute, memory, interconnect, and hardware-software tradeoffs for future GPU studies; distinguished workload-specific gains from end-to-end effects to inform engineering and product decisions.
\item Presented technical findings across architecture and product teams, developed study materials, and taught colleagues how to run and explain performance modeling workflows.
\end{itemize}

\vspace{4pt}
\role{NVIDIA}{May 2023 -- Aug 2023}{Deep Learning Performance Architect Intern}{Santa Clara, CA}
\begin{itemize}
\item Researched performance modeling of large language models on future GPU architectures.
\end{itemize}

\vspace{4pt}
\role{Micron Technology}{May 2022 -- Aug 2022}{Advanced Hardware Development Intern}{Allen, TX}
\begin{itemize}
\item Researched enhanced in-memory computing functions and memory-system design.
\end{itemize}

\section{Education}
\textbf{University of Michigan, Ann Arbor}\\
Ph.D., Computer Science and Engineering \hfill 2021 -- 2024\\
Advisor: Prof. Trevor Mudge. GPA: 4.0/4.0.\\
M.S.E., Computer Science and Engineering \hfill 2019 -- 2021\\
B.S.E., Computer Engineering \hfill 2017 -- 2019

\textbf{Shanghai Jiao Tong University}\\
B.S.E., Electrical and Computer Engineering \hfill 2015 -- 2019

\newpage
{\large\textbf{Haojie Ye}}\hfill AI Systems and GPU Architecture

\section{Selected Research}
Hardware-software co-design for recommendation inference, GPU and memory systems, graph processing, and sparse computation. Publications at HPCA, ASPLOS, MICRO, ISCA, VLDB, PACT, and other systems venues. Full publication list on \href{https://scholar.google.com/citations?user=8Mk6B6gAAAAJ}{Google Scholar}.

\paper{Palermo: Improving the Performance of Oblivious Memory using Protocol-Hardware Co-Design}{\textbf{H. Ye}, Y. Xia, Y. Chen, K.-Y. Chen, Y. Yuan, S. Deng, B. Kasikci, T. Mudge, N. Talati}{HPCA 2025}{https://doi.org/10.1109/HPCA61900.2025.00038}
Protocol and hardware co-design to overlap oblivious-memory operations while preserving security guarantees.\par\vspace{3pt}

\paper{GRACE: A Scalable Graph-Based Approach to Accelerating Recommendation Model Inference}{\textbf{H. Ye}, S. Vedula, Y. Chen, Y. Yang, A. Bronstein, R. Dreslinski, T. Mudge, N. Talati}{ASPLOS 2023}{https://doi.org/10.1145/3582016.3582029}
Graph clustering and embedding partial-sum caching to reduce recommendation inference memory traffic.\par\vspace{3pt}

\paper{Understanding the Performance and Estimating the Cost of LLM Fine-Tuning}{Y. Xia, J. Kim, Y. Chen, \textbf{H. Ye}, S. Kundu, C. Hao, N. Talati}{IISWC 2024}{https://doi.org/10.1109/IISWC63097.2024.00027}

\paper{Locality-aware Optimizations for Improving Remote Memory Latency in Multi-GPU Systems}{L. Belayneh, \textbf{H. Ye}, K.-Y. Chen, D. Blaauw, T. Mudge, R. Dreslinski, N. Talati}{PACT 2022}{https://doi.org/10.1145/3559009.3569649}

\paper{Everest: GPU-Accelerated System For Mining Temporal Motifs}{Y. Yuan, \textbf{H. Ye}, S. Vedula, W. Kaza, N. Talati}{PVLDB 17(2), 2023 / VLDB 2024}{https://doi.org/10.14778/3626292.3626299}

\paper{NDMiner: Accelerating Graph Pattern Mining Using Near Data Processing}{N. Talati, \textbf{H. Ye}, Y. Yang, L. Belayneh, K.-Y. Chen, D. Blaauw, T. Mudge, R. Dreslinski}{ISCA 2022}{https://doi.org/10.1145/3470496.3527437}

\section{Technical Skills}
\textbf{Languages:} Python, C/C++, CUDA, Verilog, SystemVerilog.\\
\textbf{Performance and architecture:} Workload characterization, performance modeling, bottleneck analysis, GPU and memory systems, parallel computing, hardware-software co-design, agent-assisted analysis.\\
\textbf{Tools:} Git; gem5, gem5-gpu, DRAMSim, Ramulator, CACTI, zsim; Cadence, Design Compiler.

\section{Teaching}
\textbf{University of Michigan --- Graduate Student Instructor}\hfill Sep 2020 -- Jan 2021\\
EECS 598, Applied Parallel Programming with GPUs. Prepared lecture slides, homework, and exams on GPU microarchitecture and CUDA programming with Prof. Reetuparna Das.

\section{Honors}
Outstanding Undergraduate Research Award (2019); James B. Angell Scholar (2019).

\end{document}
